\chapter*{循迹之外：力学的参考系变换与分析力学认知——以机械臂运动为例}
\addcontentsline{toc}{chapter}{循迹之外：力学的参考系变换与分析力学认知——以机械臂运动为例}
\addtocontents{toc}{\protect\setcounter{tocdepth}{0}}

% 保存原始的编号格式
\let\oldthesection\thesection
\let\oldthesubsection\thesubsection

% 重置section计数器并修改编号格式
\setcounter{section}{0}
\renewcommand{\thesection}{\arabic{section}}
\renewcommand{\thesubsection}{\thesection.\arabic{subsection}}

在以往的力学分析中，我们习惯于在统一的参考系下求解，这保证了矢量运算的合法和正确性。
但是，让我们考虑下面一个情况：如果一个小球在一个轻杆上从一个端点开始以匀速向另一个端点运动，
与此同时轻杆绕着这个点匀速转动，这时小球运动的切向和径向加速度又是多少？

我们直接计算会看到非常麻烦，但很明显，如果我们将小球的运动控制在轻杆的坐标系中，它将会做非常简单的匀速直线运动。
所以，有没有一个方法能够实现坐标系运动的转换？
\section{旋转矩阵和变换矩阵}
为了解决这个问题，我们需要首先对不同的坐标系构建一个联系。为了方便讨论，我们首先定义空间中矢量的表示方式。

\begin{singlefigure}[向量示意图]{向量表示}[0.45]
\end{singlefigure}

对于如上图所示的向量我们已经知道写成向量形式该如何表示，我们定义以下符号：
\[^A D = \begin{pmatrix} d_x \\ d_y \\ d_z \end{pmatrix}\]

这个符号表示：在A坐标系下，向量D可以被这样一个3*1矩阵表示。注意，对于任意一个向量，我们所在的参考系要标注于向量字母的左上角，这非常重要，良好的符号标记将极大地减轻计算的压力。

定义矢量之后，我们要特殊提到矢量的叉乘。在数学中我们一般采用公式或行列式求解，但在机器人运动学中，叉乘也等价于矩阵的乘积，例如对于一个

现在，还记得我们的目的是在不同坐标系下表示一个运动，那么不同坐标系之间的转换就是一个必要条件。首先我们考虑坐标系是相对静止
的情况，如下图所示：
\begin{singlefigure}[坐标系相对关系]{坐标系相对关系}[0.45]
\end{singlefigure}

在这个过程中我们考虑从A到B的坐标系变换，明显地，这可以被分为旋转和平移两部分。平移部分我们很容易知道可以用一个三维矢量表示这种变换，那么旋转部分呢？就得用到我们的旋转矩阵。

旋转矩阵是一个3*3的矩阵，属于SO(3)群。旋转矩阵满足以下关系：
\[^A_B R \in SO(3) = 
\left\{
\begin{pmatrix}
    r_1 & r_2 & r_3
\end{pmatrix}
\in {R}^{3\times3} \mid r_i^T r_j = \delta_{ij},r_1\times r_2 = r_3
\right\}
\]

这里的角标表示，这是在A空间作为参考系下观察B空间的姿态。角标位置同样会事关计算和理解，不同的角标也对应不同的观察关系。。

旋转矩阵与空间中的姿态方向唯一对应。很明显地，我们可以发现旋转操作是不改变矢量的长度，只是改变了其在空间中的表示，因此旋转矩阵是一个行列式为1的矩阵，且其对应的算子为等距同构。与此同时，旋转矩阵满足以下命题：
\[\forall R\in SO(3),\forall P,Q\in R^3,R(P\times Q) = RP\times RQ\]

现在我们加上A坐标系下B的位置矢量$^A O_B$，整体可以表示为齐次变换矩阵，这是一组4*4的矩阵，属于SE(3)群。这个集合中的元素均可表示为：
\[^A_B T = 
\begin{pmatrix}
    ^A_B R & ^A O_B \\
    0 & 1
\end{pmatrix}
\]

不同位姿与SE(3)群中的元素是一一对应的，这里的“位姿”同时表征了坐标系之间的旋转和位移关系。

现在，我们再考虑这样一件事：在A坐标系中变换到B坐标系，和在B坐标系中变换到A坐标系，两个过程是等价的。也就是说这是一组逆变换关系，因此，对应的旋转矩阵和齐次变换矩阵也是一组逆矩阵，即：
\[^A_B R = {}^B_A R^{-1}\qquad ^A_B T = {}^B_A T^{-1}\]

剩下的结论我们直接给出，读者自证不难：
\[\begin{pmatrix}
    ^A P\\
    1
\end{pmatrix} = {}^A_B T \begin{pmatrix}
    ^B P\\
    1
\end{pmatrix}\qquad ^B O_A = -{}^B_A R {}^A O_B
\]
\[^A_C R = {}^A_B R {}^B_C R\qquad ^A_C T = {}^A_B T {}^B_C T\]

我们现在知道旋转矩阵和变换矩阵了，那么对于一个具体的旋转该如何表示？我们先考虑一个以坐标轴为旋转轴的情况：
\begin{singlefigure}[姿态的表示]{姿态的欧拉角表示}[0.5]
\end{singlefigure}

如上图所示，假设B坐标系是在A坐标系基础上逆时针旋转$\theta$角得到，那么我们可以得到：
\[\hat{Y}_B = \cos\theta \hat{Y}_A + \sin\theta \hat{Z}_A\]
\[\hat{Z}_B = -\sin\theta \hat{Y}_A + \cos\theta \hat{Z}_A\]

因此，围绕$x$轴旋转的旋转矩阵可以表示如下：
\[{}^A_B R = R_x(\theta) = 
\begin{pmatrix}
    1 & 0 & 0 \\
    0 & \cos\theta & -\sin\theta \\
    0 & \sin\theta & \cos\theta
\end{pmatrix}
\]

同理可得绕$y$轴和$z$轴的旋转矩阵，结果如下：
\[R_y(\theta) = 
\begin{pmatrix}
    \cos\theta & 0 & \sin\theta \\
    0 & 1 & 0 \\
    -\sin\theta & 0 & \cos\theta
\end{pmatrix}
\qquad
R_z(\theta) = 
\begin{pmatrix}
    \cos\theta & -\sin\theta & 0 \\
    \sin\theta & \cos\theta & 0 \\
    0 & 0 & 1
\end{pmatrix}
\]

进一步考虑空间中的任意姿态，我们发现只要依次绕着坐标轴的三个方向进行旋转即可得到，这种过程对应变换的复合，即矩阵的乘法。相应的角包括$z-y-x$等，对应的姿态表示为$R_{z'y'x'}$，表示依次以$z$、$y$、$x$轴的顺序旋转。

但是变换中我们发现，世界坐标系和研究对象坐标系本身是两套坐标系，我们规定与研究对象结合的为联体坐标系，表示对象旋转时坐标系跟随旋转（比如设定飞机机头为$x$轴的坐标系），对应的角度称为欧拉角。而另一个绝对的坐标系为基坐标系，对应的角度称为固定角。对于同一个姿态，二者的矩阵乘法顺序相反，即：
\[R_{xyz} = R_{z'y'x'} = R_z(\alpha) R_y(\beta) R_x(\gamma)\]

其中，第一项表示操作是相对于固定坐标系的$x-y-z$操作顺序，第二项表示操作是相对于联体坐标系的$z-y-x$操作顺序。简记为“右乘联体左乘基”。

姿态可以通过三次旋转决定，那一次旋转是否可行？欧拉旋转定理指出：若刚体从初姿态作任意定点转动后呈终姿态，则必可找到一个过该点的轴$K$及角度$\theta$，则该轴为旋转轴，该轴旋转$\theta$后，刚体呈终姿态。假设轴为一个单位向量${}^A K = (k_x,k_y,k_z)^T$（注意以下公式必须要求是单位向量），由于推导比较复杂，我们直接给出这种等效轴角方法求得的旋转矩阵\footnote{在机器人运动学中，为了便于书写，$\sin \theta$和$\cos \theta$分别简写为$s\theta$和$c\theta$，$v\theta = 1-c\theta$}：
\[{}^A_B R = R_K(\theta)
\begin{pmatrix}
    k_x^2 v\theta+c\theta & k_x k_y v\theta-k_z s\theta & k_x k_z v\theta+k_y s\theta \\
    k_y k_x v\theta+k_z s\theta & k_y^2 v\theta+c\theta & k_y k_z v\theta-k_x s\theta \\
    k_z k_x v\theta-k_y s\theta & k_z k_y v\theta+k_x s\theta & k_z^2 v\theta+c\theta
\end{pmatrix}
\]
\section{微分运动学}
同上一节，我们首先定义一些符号：
\[v_c = {}^U V_{CORG}\qquad \omega_c = {}^U \omega_{CORG}\]

上面这两个符号分别表示在世界坐标系{U}下观测的z坐标系{C}的原点的速度和角速度，注意，这与${}^C v_c$是不同的，后者表示在C坐标系下观测的速度。另外，上标相同时通常会省略，例如${}^A({}^A V_Q) = {}^A V_Q$。

\begin{minipage}{0.5\textwidth}
    \begin{singlefigure}[速度微分]{速度微分}[0.95]
    \end{singlefigure}
\end{minipage}
\begin{minipage}{0.4\textwidth}
    \begin{singlefigure}[角速度微分]{角速度微分}[0.9]
    \end{singlefigure}
\end{minipage}
\\[1em]

现在，让我们来推导速度和加速度表达式。首先是速度，我们要求${}^A V_Q$和${}^B V_Q$的关系：
\[
\begin{aligned}
    {}^A V_Q &= \lim_{\Delta t\rightarrow 0}\dfrac{{}^A Q(t + \Delta t) - {}^A Q(t)}{\Delta t}\\
    &= \lim_{\Delta t\rightarrow 0}\dfrac{{}^A O_B(t + \Delta t) + {}^A_B R(t + \Delta t){}^B Q(t + \Delta t) - {}^A O_B(t) - {}^A_B R(t){}^B Q(t)}{\Delta t}\\
    &= \lim_{\Delta t\rightarrow 0}{}^A V_{BORG} + \dfrac{d}{dt}({}^A_B R {}^B Q)\\
    &= {}^A V_{BORG} + {}^A_B \dot{R} {}^B Q+ {}^A_B R {}^B V_Q\\
\end{aligned}
\]

现在，我们的问题变成了：这个旋转矩阵的求导该如何得到？从上一节最后我们知道，任意一个旋转矩阵等价于绕一个空间中的轴旋转一次得到。因此我们对等效轴角表示的旋转矩阵进行求导，首先定义角速度向量和瞬轴的归一化向量：
\[{}^A K = (k_x, k_y, k_z)^T = (\Omega_x/\dot{\phi},\Omega_y/\dot{\phi},\Omega_z/\dot{\phi})^T \qquad {}^A \Omega_B = (\Omega_x, \Omega_y, \Omega_z)^T\]

其中标量$\dot{\phi}$表示旋转速度，这是合理的因为针对瞬轴这个情况下确实是在旋转。由于在这个时间间隔内，物体是在绕这个瞬轴旋转得到$t+\Delta t$时刻的旋转矩阵，因此${}^A_B R(t+\Delta t) = {}^A R_K(\phi){}^A_B R$，故：
\[\begin{aligned}
    {}^A_B \dot{R} &= \lim_{\Delta t\rightarrow 0}\dfrac{R_K(\phi)-R_K(0)}{\Delta t}{}^A_B R \\
    &=\lim_{\phi\rightarrow 0}\dfrac{R_K(\phi) - I_3}{\phi}\dot{\phi} {}^A_B R
\end{aligned}
\]

代入$\phi\rightarrow 0$到我们在上一节的表达式，可得：
\[{}^A_B \dot{R} = \dot{\phi}\begin{pmatrix}
    0 & -k_z & k_y\\
    k_z & 0 & -k_x\\
    -k_y & k_x & 0
\end{pmatrix}{}^A_B R
= \begin{pmatrix}
    0 & -\Omega_z & \Omega_y\\
    \Omega_z & 0 & -\Omega_x\\
    -\Omega_y & \Omega_x & 0
\end{pmatrix}{}^A_B R
\]

我们把这里的前一个矩阵定义为${}^A_B S$，这是一个反对称矩阵。其中的元素都是角速度分量$\Omega_i$，因此是${}^A \Omega_B$对应的反对称矩阵。根据计算关系，我们可以得到\mgnote{事实上，矩阵的叉乘也可表示为前者对应的反对称矩阵与后者的乘积}：
\[{}^A_B S {}^A_B R = {}^A \Omega_B\times {}^A_B R\]

因此我们可以得到：
\[{}^A V_Q = {}^A V_{BORG} + {}^A_B R {}^B V_Q + {}^A\Omega_B\times {}^A_B R {}^B Q\]

同理，对于角速度递推关系，我们从${}^A_C R = {}^A_B R{}^B_C R$开始推导，利用上述反对称矩阵关系（过程略）可以得到：
\[{}^A \Omega_C = {}^A\Omega_B + {}^A_B R {}^B \Omega_C\]

同时，我们这种速度和角速度的结果可以放在不同坐标系下观察，这时只需要一个旋转矩阵把对应矢量在空间中的位姿进行旋转即可。

从前面的推导中，我们不仅能够在不同坐标系下描述一个物体的运动位置和速度、角速度，根据前面的推导，我们也发现一个旋转矩阵的导数可以用它的旋转角速度表示为一种叉乘的形式，这样我们可以很顺利地推导出想要的结果。

% 在速度层面，以上结果可以进一步描述为一种向外递推的结构，而对于力的传递，我们会假定给末端一个假想的力$f = (f_x,f_y,f_z)^T$，通过瞬时的力和力矩平衡条件倒推，属于向内迭代的过程。但是由于这种推导对应的公式需要绑定特定的坐标系建模，因此我们暂不讨论。读者感兴趣可自行尝试推导。

\section{坐标系的跳跃——我们最初的问题}
现在，让我们拿着这个工具具体讨论最初我们想解决的问题。这是一个经典的阿基米德螺线，我们发现在极坐标系下描述这个运动是很简单的，它有着径向向外的匀速和世界坐标系下的角速度恒定。

\begin{center}
    \begin{tikzpicture}[>=latex, scale=1]
        % 1. 固定坐标系 {A}
        \draw[->, thin, gray] (-3.5,0) -- (5,0) node[below] {$x_A$};
        \draw[->, thin, gray] (0,-3) -- (0,4.5) node[left] {$y_A$};
        \draw[fill] (0,0) circle (0.05) node[below left] {$O$};

        % 2. 旋转的杆（动坐标系 {B} 的轴线）
        \def\ang{40} % 当前角度
        \draw[thick, brown] (0,0) -- (\ang:5.5) node[above right] {杆 (轴 $\hat{r}$)};
        
        % 3. 角度标注
        \draw[->, thin, cyan] (1.0,0) arc (0:\ang:1.0) node[pos=0.6, right] {$\theta=\omega t$};

        % 4. 轨迹（阿基米德螺线）- 从 r=1.0 开始，让径向增长明显且图形更开阔
        \def\rstart{1.0}
        \def\vrad{1.5} % 增大径向速度系数，让螺线展开更明显
        % 角度参数 \t 从 0 到 \ang（度），绘制平滑曲线
        \draw[thick, magenta, dotted, domain=0:\ang, samples=50, smooth, variable=\t]
            plot ({(\rstart + \vrad*(\t/57.2958))*cos(\t)},
                 {(\rstart + \vrad*(\t/57.2958))*sin(\t)});

        % 5. 当前质点位置 P - 精确计算在螺线上的位置
        \coordinate (P) at (\ang:{\rstart + \vrad*(\ang/57.2958)});
        \draw[fill=black] (P) circle (0.08) node[above right] {$P$};

        % 6. 相对速度（沿杆滑动）—— 过程描述，不是结果计算
        \draw[->, red!80!black, thick] (P) -- ++(\ang:1.2) node[above] {$\mathbf{v}_r$};

        % 7. 旋转方向示意（在圆心附近画个弯箭头）
        \draw[->, thick, blue!50] (0.4,0) arc (0:80:0.4) node[pos=0.5, right] {$\omega$};
        
        % 添加辅助圆来对比显示径向运动
        \draw[dashed, gray!40] (0,0) circle (\rstart);
    \end{tikzpicture}
\end{center}

对于一个单层的模型，似乎我们可以直接进行坐标分解为切向和径向，这并不复杂，由于：
\[\dfrac{\dif}{\dif t}\hat{e}_r = \dot{\theta}\hat{e}_{\theta}\qquad \dfrac{\dif}{\dif t}\hat{e}_\theta = -\dot{\theta}\hat{e}_r\]

所以：
\[p = r\hat{e}_r\Rightarrow v = \dot{r}\hat{e}_r + r\dot{\theta}\hat{e}_{\theta}\Rightarrow a = (\ddot{r}-r\dot{\theta}^2)\hat{e}_r + (r\ddot{\theta} + 2\dot{r}\dot{\theta})\hat{e}_{\theta}\]

如果用之前的思路，我们首先规定世界坐标系为A，与此同时有一个连体坐标系与这个杆捆绑，原点与A原点重合，杆是坐标系的X轴。这时候的角速度和线速度可以很简单地表示。因此：
\[{}^A \Omega_B = (0,0,\dot{\theta})^T\qquad {}^B V_Q = (\dot{r},0,0)^T\]
\[{}^A_B R = \left(\begin{matrix}
    \cos\theta & -\sin\theta & 0\\
    \sin\theta & \cos\theta & 0\\
    0 & 0 & 1
\end{matrix}\right)\]

根据我们在上一节得到的速度表达式以及这个过程的思路，我们继续求导得到加速度的表达式，明显第一项坐标系的相对速度为0，所以：
\[\begin{aligned}
    {}^A a_Q &= \dfrac{\dif}{\dif t}\left({}^A_B R {}^B V_Q + {}^A\Omega_B\times {}^A_B R {}^B Q\right) \\
    &= {}^A_B \dot{R} {}^B V_Q + {}^A_B R {}^B \dot{V}_Q + {}^A\dot{\Omega}_B\times {}^A_B R {}^B Q + {}^A\Omega_B\times \left({}^A_B \dot{R} {}^B Q + {}^A_B R {}^B \dot{Q}\right)
\end{aligned}\]

其中 \( {}^A_B \dot{R} = {}^A\Omega_B\times {}^A_B R \)，代入计算：

第一项（牵连加速度的一部分）：
\[{}^A_B \dot{R} {}^B V_Q = ({}^A\Omega_B\times {}^A_B R)(\dot{r},0,0)^T = \dot{\theta}\hat{k}\times(\dot{r}\cos\theta,\dot{r}\sin\theta,0)^T = (-\dot{r}\dot{\theta}\sin\theta,\dot{r}\dot{\theta}\cos\theta,0)^T\]

第二项（相对加速度）：
\[{}^A_B R {}^B \dot{V}_Q = \left(\begin{matrix}
    \cos\theta & -\sin\theta & 0\\
    \sin\theta & \cos\theta & 0\\
    0 & 0 & 1
\end{matrix}\right)\left(\begin{matrix}
    \ddot{r}\\ 0\\ 0
\end{matrix}\right) = (\ddot{r}\cos\theta,\ddot{r}\sin\theta,0)^T\]

第三项（角加速度项），因为旋转是匀速的，所以这部分为0。

第四项（科里奥利加速度和向心加速度）：
\[\begin{aligned}
    {}^A\Omega_B\times \left({}^A_B \dot{R} {}^B Q + {}^A_B R {}^B \dot{Q}\right) &= \dot{\theta}\hat{k}\times\left[({}^A\Omega_B\times {}^A_B R)(r,0,0)^T + {}^A_B R(\dot{r},0,0)^T\right] \\
    &= \dot{\theta}\hat{k}\times\left[(-r\dot{\theta}\sin\theta,r\dot{\theta}\cos\theta,0)^T + (\dot{r}\cos\theta,\dot{r}\sin\theta,0)^T\right] \\
    &= \dot{\theta}\hat{k}\times(-r\dot{\theta}\sin\theta+\dot{r}\cos\theta,r\dot{\theta}\cos\theta+\dot{r}\sin\theta,0)^T \\
    &= (-r\dot{\theta}^2\cos\theta-\dot{r}\dot{\theta}\sin\theta,-r\dot{\theta}^2\sin\theta+\dot{r}\dot{\theta}\cos\theta,0)^T
\end{aligned}\]

将所有项相加，得到世界坐标系下的加速度：
\[\begin{aligned}
    {}^A a_Q &= (\ddot{r}\cos\theta-r\dot{\theta}^2\cos\theta-2\dot{r}\dot{\theta}\sin\theta,\ddot{r}\sin\theta-r\dot{\theta}^2\sin\theta+2\dot{r}\dot{\theta}\cos\theta,0)^T \\
    &= (\ddot{r}-r\dot{\theta}^2)(\cos\theta,\sin\theta,0)^T + (2\dot{r}\dot{\theta})(-\sin\theta,\cos\theta,0)^T
\end{aligned}\]

这正是我们在极坐标系下得到的结果：径向加速度 \( (\ddot{r}-r\dot{\theta}^2)\hat{e}_r \) 和切向加速度 \( (r\ddot{\theta}+2\dot{r}\dot{\theta})\hat{e}_{\theta} \)（本题中 \( \ddot{\theta}=0 \)）。

对比上面的过程，你可能会认为这个旋转矩阵的计算是很冗余的。甚至来说，我自然可以表示出这个点的绝对位置坐标然后直接求导，得到的结果与上面的矩阵表示是完全一致的。但是，旋转矩阵带来的坐标系转换是一种通用的形式，当坐标系之间的关系变得越来越复杂的时候，这种递推关系式将保证结果是可达可求的，只要给到必要的信息确保这个点是唯一的。

那么，我们进一步讨论一个更为复杂的结果，例如一个空间中的杆，它绕着Z轴旋转，杆自己是联体坐标系X轴；在这个端点处继续有一个杆，绕着这个联体坐标系的X轴旋转，再继续在这个坐标系下有一个杆绕着第二个联体坐标系的X轴旋转；最后再在第三个联体坐标系的Z轴方向给一个平动速度。如下图所示：
\begin{singlefigure}[一种复杂运动的例子]{复杂牵连运动}[0.5]
\end{singlefigure}

很明显，这个时候单独在每个连杆中观察下一个杆的运动是极为便捷的，但是在世界坐标系下观察则很难直接写出这个绝对位置并顺利求导——即使写出来也是很长的一串，如果加上平动而非转动关节就会更麻烦。

事实上，结合我们在上面推到的结论以及这些物理量与在世界坐标系下的观察关系：${}^i v_i = {}^i_U R {}^U V_i$和${}^i \omega_i = {}^i_U R {}^U \Omega_i$（这里表示的含义是在坐标系i下观察的世界坐标系下的物体的运动），假定每个关节的旋转轴或者平动方向都在Z轴方向，那么：
\[{}^{i+1} v_{i+1} = {}^{i+1}_i R ({}^i v_i + {}^i \omega_i \times {}^i P_{i+1}) + \dot{d}_{i+1}{}^{i+1}\hat{Z}_{i+1}\]
\[{}^{i+1} \omega_{i+1} = {}^{i+1}_i R {}^i \omega_i + \dot{\theta}_{i+1} {}^{i+1}\hat{Z}_{i+1}\]

根据以上迭代关系不同坐标系下对同一矢量的观察关系$v_N = {}^0_N R {}^N v_N$和$\omega_N = {}^0_N R {}^N \omega_N$即可求出世界坐标系下的运动——只要我们能够顺利表示出这里的旋转矩阵位置关系，例如常用的有非标准D-H法等，此处不赘述。很显然，在世界坐标系U下，迭代的初始值对应的速度和角速度都是0。



% 恢复原始的编号格式
\let\thesection\oldthesection
\let\thesubsection\oldthesubsection
\addtocontents{toc}{\protect\setcounter{tocdepth}{2}}